\documentclass[12pt,a4paper]{article}

% --- CẤU HÌNH TIẾNG VIỆT ---
\usepackage[T5]{fontenc}
\usepackage[utf8]{inputenc}
\usepackage[vietnam]{babel}
\usepackage{times}

\usepackage{amsmath}
\usepackage{graphicx}
\usepackage{float}
\usepackage{geometry}
\usepackage{hyperref}
\usepackage{listings}
\usepackage{color}
\usepackage{xcolor}
\usepackage{array}
\usepackage{booktabs}
\usepackage{tabularx}
\usepackage{multirow}
\usepackage{fancyhdr}
\usepackage{titlesec}
\usepackage{enumitem}
\usepackage{mdframed}
\usepackage{colortbl}

\geometry{
  a4paper,
  left=30mm,
  right=20mm,
  top=25mm,
  bottom=25mm
}

% --- Cấu hình header/footer ---
\pagestyle{fancy}
\fancyhf{}
\fancyhead[L]{\small \textit{Báo cáo thực tập – SmartElevator Camera}}
\fancyhead[R]{\small \textit{VNPT Media}}
\fancyfoot[C]{\thepage}
\renewcommand{\headrulewidth}{0.4pt}

% --- Màu sắc ---
\definecolor{vnptblue}{RGB}{0, 82, 156}
\definecolor{lightblue}{RGB}{220, 235, 252}
\definecolor{darkgray}{RGB}{64, 64, 64}
\definecolor{weekbg}{RGB}{240, 248, 255}
\definecolor{resultbg}{RGB}{240, 255, 240}
\definecolor{codegreen}{rgb}{0,0.6,0}
\definecolor{codegray}{rgb}{0.5,0.5,0.5}
\definecolor{codepurple}{rgb}{0.58,0,0.82}
\definecolor{backcolour}{rgb}{0.95,0.95,0.92}

% --- Cấu hình code listing ---
\lstdefinestyle{cppstyle}{
    backgroundcolor=\color{backcolour},
    commentstyle=\color{codegreen},
    keywordstyle=\color{blue}\bfseries,
    numberstyle=\tiny\color{codegray},
    stringstyle=\color{codepurple},
    basicstyle=\ttfamily\footnotesize,
    breakatwhitespace=false,
    breaklines=true,
    captionpos=b,
    keepspaces=true,
    numbers=left,
    numbersep=5pt,
    showspaces=false,
    showstringspaces=false,
    showtabs=false,
    tabsize=2,
    frame=single
}
\lstset{style=cppstyle}

% --- Cấu hình tiêu đề section ---
\titleformat{\section}
  {\Large\bfseries\color{vnptblue}}
  {\thesection.}{0.5em}{}[\titlerule]

\titleformat{\subsection}
  {\large\bfseries\color{darkgray}}
  {\thesubsection.}{0.5em}{}

% --- Hyperref ---
\hypersetup{
    colorlinks=true,
    linkcolor=vnptblue,
    urlcolor=blue,
    pdftitle={Báo cáo thực tập hàng tuần - SmartElevator},
}

% --- Môi trường khung tuần ---
\newmdenv[
  linecolor=vnptblue,
  linewidth=1.5pt,
  backgroundcolor=weekbg,
  innerleftmargin=10pt,
  innerrightmargin=10pt,
  innertopmargin=8pt,
  innerbottommargin=8pt
]{weekframe}

\begin{document}

% =======================================================
% TRANG BÌA
% =======================================================
\begin{titlepage}
    \centering
    \vspace*{0.5cm}

    {\large \textbf{CỘNG HÒA XÃ HỘI CHỦ NGHĨA VIỆT NAM}} \\
    {\large \textbf{Độc lập -- Tự do -- Hạnh phúc}} \\
    \vspace{0.3cm}
    \rule{8cm}{1.2pt}

    \vspace{1.5cm}

    {\large \textbf{ĐƠN VỊ THỰC TẬP: VNPT MEDIA}} \\[0.3cm]
    {\large Phòng Nghiên cứu và Phát triển Công nghệ} \\

    \vspace{2cm}

    \rule{14cm}{1pt} \\[0.5cm]
    {\Huge \bfseries \color{vnptblue} BÁO CÁO THỰC TẬP} \\[0.4cm]
    {\LARGE \bfseries HÀNG TUẦN} \\[0.3cm]
    \rule{14cm}{1pt}

    \vspace{1.5cm}

    {\Large \bfseries Dự án: Hệ thống Thang máy Thông minh} \\[0.3cm]
    {\large \textit{SmartElevator Camera -- Nhận diện khuôn mặt AI}} \\[0.3cm]
    {\normalsize Ngôn ngữ: C++17 \quad | \quad Thư viện: OpenCV, SQLite3, Win32 API}

    \vspace{2.5cm}

    \begin{table}[H]
        \centering
        \large
        \begin{tabular}{ll}
            \textbf{Sinh viên thực tập} & : Nguyễn Đức Tuấn \\[0.3cm]
            \textbf{Mentor hướng dẫn}   & : Kỹ sư VNPT Media \\[0.3cm]
            \textbf{Thời gian báo cáo}  & : Tuần 1 \& Tuần 2 \\[0.3cm]
            \textbf{Năm thực tập}       & : 2026 \\
        \end{tabular}
    \end{table}

    \vfill
    \rule{8cm}{0.5pt} \\[0.3cm]
    {\textbf{Hà Nội, tháng 7 năm 2026}}
\end{titlepage}

\newpage
\pagenumbering{roman}
\tableofcontents
\newpage
\pagenumbering{arabic}

% =======================================================
% THÔNG TIN CHUNG
% =======================================================
\section*{THÔNG TIN CHUNG VỀ KỲ THỰC TẬP}
\addcontentsline{toc}{section}{Thông tin chung về kỳ thực tập}

\begin{table}[H]
    \centering
    \renewcommand{\arraystretch}{1.4}
    \begin{tabular}{|>{\bfseries}p{5cm}|p{9.5cm}|}
        \hline
        \rowcolor{lightblue}
        \multicolumn{2}{|c|}{\bfseries THÔNG TIN SINH VIÊN THỰC TẬP} \\
        \hline
        Họ và tên sinh viên   & Nguyễn Đức Tuấn \\
        \hline
        Đơn vị thực tập       & VNPT Media -- Phòng Nghiên cứu \& Phát triển Công nghệ \\
        \hline
        Tên đề tài thực tập   & Hệ thống Thang máy Thông minh điều khiển bằng Nhận diện khuôn mặt AI (SmartElevator Camera) \\
        \hline
        Ngôn ngữ lập trình    & C++ (Chuẩn C++17) \\
        \hline
        Thư viện / Công cụ    & OpenCV 4.x, SQLite3, Win32 API, CMake, MSYS2/UCRT64 \\
        \hline
        Thời gian báo cáo     & Tuần 1 và Tuần 2 của kỳ thực tập năm 2026 \\
        \hline
        Mentor hướng dẫn      & Kỹ sư phụ trách -- VNPT Media \\
        \hline
    \end{tabular}
    \caption{Thông tin tổng quan kỳ thực tập}
\end{table}

\newpage

% =======================================================
% TUẦN 1
% =======================================================
\section{BÁO CÁO TUẦN 1 -- NGHIÊN CỨU VÀ THIẾT KẾ HỆ THỐNG}

\subsection{Tổng quan mục tiêu tuần 1}

\begin{weekframe}
\textbf{Mục tiêu tuần 1:} Tiếp nhận đề tài, nghiên cứu tài liệu yêu cầu nghiệp vụ, khảo sát công nghệ, thiết kế cấu trúc dữ liệu và luồng xử lý hệ thống để trình Mentor duyệt trước khi bắt đầu coding ở tuần 2.
\end{weekframe}

\vspace{0.5cm}

\subsection{Nhật ký công việc chi tiết từng ngày}

\begin{table}[H]
    \centering
    \renewcommand{\arraystretch}{1.6}
    \begin{tabularx}{\textwidth}{|>{\bfseries\centering\arraybackslash}p{2cm}|>{\arraybackslash}X|>{\arraybackslash}p{4.5cm}|}
        \hline
        \rowcolor{lightblue}
        \textbf{Thời gian} & \textbf{Nội dung công việc thực hiện} & \textbf{Kết quả / Sản phẩm yêu cầu} \\
        \hline
        Thứ Hai &
        Làm thủ tục nhận việc, làm quen môi trường và tiếp nhận đề tài từ Mentor. &
        Hiểu rõ cơ cấu phòng ban và phạm vi đề tài thực tập. \\
        \hline
        Thứ Ba &
        Nghiên cứu sâu về tài liệu mô tả yêu cầu nghiệp vụ (SRS) và làm rõ các ca sử dụng (Use Case). &
        Danh sách các yêu cầu chức năng và phi chức năng của hệ thống. \\
        \hline
        Thứ Tư &
        Khảo sát công nghệ triển khai (Ngôn ngữ C++, OOP, các thư viện quản lý luồng dữ liệu). &
        Tài liệu nghiên cứu công nghệ và kiến trúc đề xuất. \\
        \hline
        Thứ Năm &
        Thiết kế cấu trúc dữ liệu và sơ đồ cơ sở dữ liệu (Database Schema) / cấu trúc phân lớp. &
        Sơ đồ quan hệ thực thể hoặc Class Diagram chi tiết. \\
        \hline
        Thứ Sáu &
        Thiết kế luồng xử lý dữ liệu và xây dựng tài liệu thiết kế hệ thống sơ bộ để trình Mentor duyệt. &
        Báo cáo thiết kế hệ thống tuần 1 được Mentor VNPT Media thông qua. \\
        \hline
    \end{tabularx}
    \caption{Nhật ký công việc Tuần 1}
    \label{tab:week1_log}
\end{table}

\subsection{Chi tiết nội dung nghiên cứu}

\subsubsection{Thứ Hai -- Tiếp nhận đề tài và làm quen môi trường}

Ngày đầu tiên của kỳ thực tập, sinh viên đã hoàn thành các thủ tục hành chính và được Mentor VNPT Media giới thiệu tổng quan về:
\begin{itemize}[leftmargin=1.5cm]
    \item Cơ cấu tổ chức của Phòng Nghiên cứu \& Phát triển Công nghệ.
    \item Quy trình làm việc, báo cáo hàng tuần và tiêu chí đánh giá.
    \item Đề tài được giao: \textbf{Xây dựng hệ thống thang máy thông minh điều khiển bằng nhận diện khuôn mặt AI (SmartElevator Camera)}.
    \item Môi trường phát triển: Windows 10/11, MSYS2/UCRT64, CMake, OpenCV 4.x, SQLite3.
\end{itemize}

\textbf{Kết quả đạt được:} Nắm rõ yêu cầu đề tài, cài đặt môi trường lập trình cơ bản và hiểu được kỳ vọng của Mentor về sản phẩm cuối kỳ.

\subsubsection{Thứ Ba -- Nghiên cứu tài liệu SRS và Use Case}

Sinh viên đã nghiên cứu kỹ lưỡng tài liệu mô tả yêu cầu nghiệp vụ (Software Requirements Specification -- SRS) và xác định các ca sử dụng chính của hệ thống.

\paragraph{Yêu cầu chức năng (Functional Requirements):}
\begin{enumerate}[leftmargin=1.5cm]
    \item \textbf{Đăng ký khuôn mặt:} Hệ thống phải cho phép đăng ký cư dân mới với thông tin: Họ tên, số căn hộ, tầng đăng ký và chụp ảnh khuôn mặt mẫu.
    \item \textbf{Nhận diện thời gian thực:} Phát hiện và nhận diện khuôn mặt từ luồng video camera, phân biệt cư dân hợp lệ và người lạ.
    \item \textbf{Điều khiển thang máy:} Tự động gửi lệnh điều khiển tầng đích qua giao thức Serial UART khi nhận diện thành công.
    \item \textbf{Quản lý cơ sở dữ liệu:} Lưu trữ thông tin cư dân bền vững bằng SQLite3.
    \item \textbf{Giao diện trực quan:} Hiển thị mô phỏng trạng thái thang máy (đang chờ, mở cửa, di chuyển, đến nơi).
\end{enumerate}

\paragraph{Yêu cầu phi chức năng (Non-Functional Requirements):}
\begin{itemize}[leftmargin=1.5cm]
    \item \textbf{Hiệu năng:} Nhận diện thời gian thực đạt tối thiểu 25 FPS.
    \item \textbf{Độ chính xác:} Tỷ lệ nhận diện đúng $\geq 90\%$ trong điều kiện ánh sáng ổn định.
    \item \textbf{Bảo mật:} Chỉ cư dân đã đăng ký mới được kích hoạt thang máy.
    \item \textbf{Khả năng mở rộng:} Hệ thống hỗ trợ thêm cư dân mới mà không cần thay đổi cấu trúc lõi.
\end{itemize}

\subsubsection{Thứ Tư -- Khảo sát công nghệ triển khai}

Sau khi xác định rõ yêu cầu, sinh viên tiến hành khảo sát và so sánh các công nghệ phù hợp:

\begin{table}[H]
    \centering
    \renewcommand{\arraystretch}{1.4}
    \begin{tabular}{|>{\bfseries}p{3cm}|p{4.5cm}|p{6cm}|}
        \hline
        \rowcolor{lightblue}
        \textbf{Công nghệ} & \textbf{Lý do lựa chọn} & \textbf{Vai trò trong hệ thống} \\
        \hline
        C++ (C++17) & Hiệu năng cao, kiểm soát bộ nhớ trực tiếp, hỗ trợ tốt các thư viện xử lý ảnh. & Ngôn ngữ lập trình chính của toàn bộ dự án. \\
        \hline
        OpenCV 4.x & Thư viện xử lý ảnh mã nguồn mở hàng đầu, tích hợp sẵn thuật toán Haar Cascade và LBPH. & Xử lý khung hình, phát hiện và nhận diện khuôn mặt. \\
        \hline
        SQLite3 & Cơ sở dữ liệu nhúng nhẹ, không cần server, phù hợp ứng dụng desktop. & Lưu trữ thông tin cư dân và quản lý dữ liệu bền vững. \\
        \hline
        Win32 API & API hệ thống Windows, hỗ trợ giao tiếp cổng COM nối tiếp (UART) trực tiếp. & Gửi lệnh điều khiển tầng đến phần cứng thang máy. \\
        \hline
        CMake & Công cụ build đa nền tảng, tích hợp tốt với MSYS2/MinGW. & Quản lý và tự động hóa quá trình biên dịch. \\
        \hline
    \end{tabular}
    \caption{Bảng so sánh và lý do lựa chọn công nghệ}
\end{table}

\subsubsection{Thứ Năm -- Thiết kế cấu trúc dữ liệu và Database Schema}

Sinh viên tiến hành thiết kế cấu trúc lớp (Class Diagram) và sơ đồ cơ sở dữ liệu cho hệ thống:

\paragraph{Cấu trúc bảng cơ sở dữ liệu SQLite:}
\begin{table}[H]
    \centering
    \renewcommand{\arraystretch}{1.4}
    \begin{tabular}{|>{\ttfamily}l|>{\ttfamily}l|p{8cm}|}
        \hline
        \rowcolor{lightblue}
        \textbf{Tên cột} & \textbf{Kiểu dữ liệu} & \textbf{Mô tả} \\
        \hline
        id & INTEGER & Khóa chính, tự động tăng, dùng làm nhãn (label) cho LBPH. \\
        \hline
        name & TEXT & Họ tên cư dân (duy nhất, không được rỗng). \\
        \hline
        apartment & TEXT & Ký hiệu căn hộ (VD: P502, P1005). \\
        \hline
        target\_floor & INTEGER & Số tầng đích đăng ký cố định của cư dân. \\
        \hline
    \end{tabular}
    \caption{Database Schema -- Bảng \texttt{residents}}
\end{table}

\paragraph{Cấu trúc lớp đề xuất:}
\begin{itemize}[leftmargin=1.5cm]
    \item \textbf{Lớp \texttt{DatabaseManager}:} Quản lý toàn bộ thao tác CRUD với SQLite3.
    \item \textbf{Struct \texttt{Resident}:} Đại diện thông tin một cư dân (id, name, apartment, targetFloor).
    \item \textbf{Struct \texttt{Elevator}:} Máy trạng thái điều khiển logic thang máy.
    \item \textbf{Hàm \texttt{enrollFace()}:} Chụp và lưu ảnh khuôn mặt mẫu cho cư dân.
    \item \textbf{Hàm \texttt{trainRecognizer()}:} Huấn luyện mô hình LBPH và lưu ra file \texttt{.yml}.
    \item \textbf{Hàm \texttt{updateElevator()}:} Cập nhật trạng thái thang máy theo từng chu kỳ thời gian.
\end{itemize}

\subsubsection{Thứ Sáu -- Thiết kế luồng xử lý và trình duyệt Mentor}

Sinh viên tổng hợp toàn bộ thiết kế thành tài liệu hệ thống sơ bộ và trình bày với Mentor VNPT Media.

\paragraph{Luồng xử lý chính được đề xuất:}
\begin{center}
    \textit{Camera} $\rightarrow$ \textit{Grayscale + EqualizeHist} $\rightarrow$ \textit{Haar Cascade phát hiện mặt} $\rightarrow$ \textit{LBPH nhận diện} $\rightarrow$ \textit{So sánh ngưỡng} $\rightarrow$ \textit{Truy vấn SQLite} $\rightarrow$ \textit{Gửi lệnh COM}
\end{center}

\textbf{Kết quả:} Tài liệu thiết kế hệ thống tuần 1 được Mentor VNPT Media xem xét và \textbf{thông qua}.

\subsection{Tổng kết tuần 1}

\begin{table}[H]
    \centering
    \renewcommand{\arraystretch}{1.4}
    \begin{tabular}{|p{7cm}|p{7cm}|}
        \hline
        \rowcolor{lightblue}
        \textbf{Công việc đã hoàn thành} & \textbf{Ghi chú / Vấn đề phát sinh} \\
        \hline
        Tiếp nhận và hiểu đề tài, xác định phạm vi. & Không có vấn đề phát sinh. \\
        \hline
        Nghiên cứu SRS, liệt kê đầy đủ yêu cầu hệ thống. & Cần làm rõ luồng xử lý chống spam lệnh khi cư dân đứng lâu. \\
        \hline
        Khảo sát và lựa chọn công nghệ (C++, OpenCV, SQLite3, Win32). & Xác nhận MSYS2/UCRT64 là môi trường build chuẩn. \\
        \hline
        Thiết kế Database Schema và Class Diagram. & Mentor yêu cầu bổ sung trường \texttt{apartment} vào bảng cư dân. \\
        \hline
        Trình bày và được Mentor duyệt tài liệu thiết kế. & Mentor phê duyệt với một số chỉnh sửa nhỏ về giao thức COM. \\
        \hline
    \end{tabular}
    \caption{Tổng kết công việc Tuần 1}
\end{table}

\newpage

% =======================================================
% TUẦN 2
% =======================================================
\section{BÁO CÁO TUẦN 2 -- TRIỂN KHAI VÀ LẬP TRÌNH HỆ THỐNG}

\subsection{Tổng quan mục tiêu tuần 2}

\begin{weekframe}
\textbf{Mục tiêu tuần 2:} Dựa trên thiết kế đã được Mentor duyệt, tiến hành triển khai toàn bộ mã nguồn C++ của hệ thống -- từ module DatabaseManager, nhận diện khuôn mặt AI (LBPH), máy trạng thái thang máy đến giao diện hiển thị và giao tiếp phần cứng qua cổng COM.
\end{weekframe}

\vspace{0.5cm}

\subsection{Nhật ký công việc chi tiết từng ngày}

\begin{table}[H]
    \centering
    \renewcommand{\arraystretch}{1.6}
    \begin{tabularx}{\textwidth}{|>{\bfseries\centering\arraybackslash}p{2cm}|>{\arraybackslash}X|>{\arraybackslash}p{4.5cm}|}
        \hline
        \rowcolor{lightblue}
        \textbf{Thời gian} & \textbf{Nội dung công việc thực hiện} & \textbf{Kết quả / Sản phẩm yêu cầu} \\
        \hline
        Thứ Hai &
        Cài đặt và cấu hình môi trường build: MSYS2/UCRT64, CMake, OpenCV 4.x, SQLite3. Viết \texttt{CMakeLists.txt} và kiểm tra build thành công. &
        Môi trường build hoạt động, sinh ra file thực thi mẫu không lỗi. \\
        \hline
        Thứ Ba &
        Lập trình module \texttt{DatabaseManager} (\texttt{database.h} / \texttt{database.cpp}): kết nối SQLite3, tạo bảng, các hàm CRUD cư dân và quản lý ảnh khuôn mặt. &
        Module Database hoàn chỉnh, test thêm/xóa/truy vấn cư dân thành công. \\
        \hline
        Thứ Tư &
        Lập trình chức năng nhận diện khuôn mặt AI: hàm \texttt{enrollFace()} chụp ảnh mẫu, \texttt{trainRecognizer()} huấn luyện mô hình LBPH và \texttt{isModelUpToDate()} kiểm tra cache model. &
        Hệ thống nhận diện được khuôn mặt, lưu và tải model \texttt{lbph\_model.yml} thành công. \\
        \hline
        Thứ Năm &
        Lập trình máy trạng thái thang máy (\texttt{Elevator State Machine}) 7 trạng thái và giao diện đồ họa mô phỏng (\texttt{drawElevatorUI()}) bằng OpenCV. &
        Giao diện thang máy hoạt động mượt: mở/đóng cửa, di chuyển giữa các tầng, hiển thị thông tin cư dân. \\
        \hline
        Thứ Sáu &
        Tích hợp toàn bộ module vào \texttt{main()}, lập trình giao tiếp cổng COM (Win32 API), kiểm thử end-to-end và debug hệ thống. &
        Hệ thống hoàn chỉnh: nhận diện mặt $\rightarrow$ gửi lệnh COM $\rightarrow$ thang máy mô phỏng chạy đúng. \\
        \hline
    \end{tabularx}
    \caption{Nhật ký công việc Tuần 2}
    \label{tab:week2_log}
\end{table}

\subsection{Chi tiết kỹ thuật từng ngày}

\subsubsection{Thứ Hai -- Cấu hình môi trường và CMakeLists.txt}

Sinh viên cài đặt và cấu hình toàn bộ môi trường phát triển:
\begin{itemize}[leftmargin=1.5cm]
    \item Cài đặt MSYS2/UCRT64, compiler GCC/G++ phiên bản mới nhất.
    \item Cài đặt thư viện OpenCV 4.x và SQLite3 qua pacman (MSYS2 package manager).
    \item Viết tệp \texttt{CMakeLists.txt} cấu hình build project:
\end{itemize}

\begin{lstlisting}[language=make, caption={Tệp CMakeLists.txt cấu hình build dự án}]
cmake_minimum_required(VERSION 3.10)
project(SmartElevatorCamera)

set(CMAKE_CXX_STANDARD 17)
set(OpenCV_DIR "C:/msys64/ucrt64/lib/cmake/opencv4")

find_package(OpenCV REQUIRED)
find_package(SQLite3 REQUIRED)

# Khai bao cac file nguon
add_executable(SmartElevatorCamera main.cpp database.cpp)

# Lien ket thu vien
target_link_libraries(SmartElevatorCamera
    ${OpenCV_LIBS}
    ${SQLite3_LIBRARIES}
)
\end{lstlisting}

\subsubsection{Thứ Ba -- Lập trình Module DatabaseManager}

Module \texttt{DatabaseManager} được xây dựng theo mô hình OOP trong hai tệp \texttt{database.h} và \texttt{database.cpp}, phụ trách toàn bộ thao tác với cơ sở dữ liệu SQLite3.

\begin{table}[H]
    \centering
    \renewcommand{\arraystretch}{1.4}
    \begin{tabular}{|>{\ttfamily}p{5.5cm}|p{9cm}|}
        \hline
        \rowcolor{lightblue}
        \textbf{Phương thức} & \textbf{Mô tả chức năng} \\
        \hline
        openDatabase(path) & Mở hoặc tạo mới tệp CSDL SQLite3 tại đường dẫn chỉ định. \\
        \hline
        createTables() & Khởi tạo bảng \texttt{residents} nếu chưa tồn tại. \\
        \hline
        addResident(name, apt, floor) & Thêm một cư dân mới vào CSDL. \\
        \hline
        deleteResident(id) & Xóa thông tin cư dân theo ID. \\
        \hline
        getAllResidents() & Trả về danh sách toàn bộ cư dân đã đăng ký. \\
        \hline
        getFloorByResidentId(id) & Lấy số tầng đăng ký của một cư dân theo ID. \\
        \hline
        saveFaceImage(id, img, idx) & Lưu ảnh mặt mẫu vào thư mục \texttt{faces/resident\_<id>/}. \\
        \hline
        loadAllFaceTrainingData() & Đọc toàn bộ ảnh JPG trong thư mục \texttt{faces/} để huấn luyện AI. \\
        \hline
    \end{tabular}
    \caption{Danh sách phương thức của lớp \texttt{DatabaseManager}}
\end{table}

\begin{lstlisting}[language=C++, caption={Khởi tạo cấu trúc bảng cư dân trong SQLite3}]
// Tao bang residents neu chua ton tai
const char* sql = R"(
    CREATE TABLE IF NOT EXISTS residents (
        id           INTEGER PRIMARY KEY AUTOINCREMENT,
        name         TEXT    NOT NULL UNIQUE,
        apartment    TEXT    NOT NULL,
        target_floor INTEGER NOT NULL
    );
)";
sqlite3_exec(db, sql, nullptr, nullptr, nullptr);
\end{lstlisting}

\subsubsection{Thứ Tư -- Lập trình nhận diện khuôn mặt AI (OpenCV LBPH)}

Đây là phần cốt lõi của hệ thống. Sinh viên hiện thực ba hàm quan trọng:

\paragraph{1. Hàm \texttt{enrollFace()} -- Thu thập ảnh khuôn mặt mẫu:}
\begin{lstlisting}[language=C++, caption={Hàm enrollFace() -- Chụp và lưu ảnh mẫu cư dân}]
void enrollFace(VideoCapture& cap, CascadeClassifier& faceCascade,
                DatabaseManager& db, const Resident& resident) {
    int sampleCount = db.getFaceImageCount(resident.id);
    Mat frame, gray;
    while (true) {
        cap >> frame;
        cvtColor(frame, gray, COLOR_BGR2GRAY);
        equalizeHist(gray, gray);

        vector<Rect> faces;
        faceCascade.detectMultiScale(gray, faces, 1.1, 5, 0, Size(80,80));

        int key = waitKey(30) & 0xFF;
        if ((key == 's' || key == 'S') && !faces.empty()) {
            Mat faceROI = gray(faces[0]);
            Mat resized;
            resize(faceROI, resized, Size(100, 100));
            // Luu anh mau vao faces/resident_<id>/face_<n>.jpg
            db.saveFaceImage(resident.id, resized, sampleCount++);
        }
        // Du 30 mau hoac nhan Q de thoat
        if (sampleCount >= ENROLLMENT_SAMPLES || key == 'q') break;
    }
}
\end{lstlisting}

\paragraph{2. Hàm \texttt{trainRecognizer()} -- Huấn luyện và lưu mô hình LBPH:}
\begin{lstlisting}[language=C++, caption={Huấn luyện mô hình LBPH và lưu cache}]
Ptr<LBPHFaceRecognizer> trainRecognizer(DatabaseManager& db,
                                         const string& modelPath) {
    // Tai toan bo anh va nhan tu thu muc faces/
    FaceTrainingData data = db.loadAllFaceTrainingData();
    auto recognizer = LBPHFaceRecognizer::create();
    recognizer->train(data.images, data.labels);

    // Luu mo hinh ra file .yml de tai su dung khi khoi dong lai
    recognizer->save(modelPath); // "lbph_model.yml"
    return recognizer;
}
\end{lstlisting}

\paragraph{3. Hàm \texttt{isModelUpToDate()} -- Kiểm tra tính hiệu lực của cache model:}
\begin{lstlisting}[language=C++, caption={So sánh thời gian file model và ảnh mẫu}]
bool isModelUpToDate(const string& modelPath, const string& facesDir) {
    LARGE_INTEGER modelTime = getFileModTime(modelPath);
    if (modelTime.QuadPart == 0) return false; // Model chua ton tai

    WIN32_FIND_DATAA ffd;
    HANDLE hFind = FindFirstFileA((facesDir+"\\*\\*.jpg").c_str(), &ffd);
    if (hFind == INVALID_HANDLE_VALUE) return true;

    do {
        LARGE_INTEGER imgTime = {0};
        imgTime.LowPart  = ffd.ftLastWriteTime.dwLowDateTime;
        imgTime.HighPart = (LONG)ffd.ftLastWriteTime.dwHighDateTime;
        if (imgTime.QuadPart > modelTime.QuadPart)
            return false; // Co anh moi -> can huan luyen lai
    } while (FindNextFileA(hFind, &ffd));
    FindClose(hFind);
    return true; // Model con hieu luc, tai truc tiep
}
\end{lstlisting}

\subsubsection{Thứ Năm -- Máy trạng thái thang máy và giao diện đồ họa}

Sinh viên thiết kế và hiện thực \textbf{State Machine} cho thang máy với 7 trạng thái hoạt động:

\begin{table}[H]
    \centering
    \renewcommand{\arraystretch}{1.4}
    \begin{tabular}{|>{\ttfamily}p{5.5cm}|p{9cm}|}
        \hline
        \rowcolor{lightblue}
        \textbf{Trạng thái} & \textbf{Mô tả hành vi} \\
        \hline
        ELEVATOR\_IDLE & Thang máy chờ, camera quét khuôn mặt liên tục. \\
        \hline
        ELEVATOR\_DOOR\_OPENING\_BEFORE & Cửa đang mở từ từ trong 1000ms khi nhận diện thành công. \\
        \hline
        ELEVATOR\_BOARDING & Cửa mở hoàn toàn, chờ cư dân bước vào trong 2000ms. \\
        \hline
        ELEVATOR\_DOOR\_CLOSING\_BEFORE & Cửa đang đóng từ từ trước khi di chuyển (1000ms). \\
        \hline
        ELEVATOR\_MOVING & Thang máy đang di chuyển, mỗi tầng mất 1500ms mô phỏng. \\
        \hline
        ELEVATOR\_ARRIVED & Đã đến tầng đích, cửa mở và chờ cư dân ra trong 3000ms. \\
        \hline
        ELEVATOR\_DOOR\_CLOSING\_AFTER & Cửa đóng sau khi cư dân ra, quay về trạng thái IDLE. \\
        \hline
    \end{tabular}
    \caption{Bảng mô tả các trạng thái của máy trạng thái thang máy}
\end{table}

Hàm \texttt{drawElevatorUI()} được lập trình để vẽ giao diện mô phỏng thang máy bằng OpenCV bao gồm:
\begin{itemize}[leftmargin=1.5cm]
    \item \textbf{LED Panel hiển thị tầng} với màu xanh/cam nhấp nháy khi di chuyển.
    \item \textbf{Thông tin cư dân}: Tên, số căn hộ, tầng đích, trạng thái hiện tại.
    \item \textbf{Cabin thang máy}: Mô phỏng cửa trượt hai bên mở/đóng theo animation.
    \item \textbf{Thanh vị trí trục tầng}: Hiển thị vị trí cabin theo thời gian thực.
    \item \textbf{Ảnh khuôn mặt cư dân} được nhúng trực tiếp bên trong cabin.
\end{itemize}

\subsubsection{Thứ Sáu -- Tích hợp, giao tiếp COM và kiểm thử end-to-end}

\paragraph{Lập trình giao tiếp cổng COM -- Win32 API:}
\begin{lstlisting}[language=C++, caption={Khởi tạo cổng COM3 và hàm gửi lệnh tầng}]
// Khoi tao cong COM3 voi cau hinh 9600 baud, 8N1
HANDLE hSerial = CreateFileA("\\\\.\\COM3",
    GENERIC_READ | GENERIC_WRITE, 0, NULL,
    OPEN_EXISTING, FILE_ATTRIBUTE_NORMAL, NULL);

DCB dcbSerialParams = {0};
dcbSerialParams.BaudRate = CBR_9600;
dcbSerialParams.ByteSize = 8;
dcbSerialParams.StopBits = ONESTOPBIT;
dcbSerialParams.Parity   = NOPARITY;
SetCommState(hSerial, &dcbSerialParams);

// Gui 1 byte ma tang qua cong COM
// Vi du: Tang 5 -> gui ky tu ASCII '5' (0x35)
void sendData(HANDLE hSerial, char data) {
    DWORD bytesWritten;
    WriteFile(hSerial, &data, 1, &bytesWritten, NULL);
}
\end{lstlisting}

\paragraph{Tích hợp trong vòng lặp chính \texttt{main()}:}

Logic chính của chương trình trong vòng lặp nhận diện thời gian thực:
\begin{lstlisting}[language=C++, caption={Vòng lặp chính -- nhận diện và điều khiển thang máy}]
while (true) {
    cap >> frame;
    cvtColor(frame, gray, COLOR_BGR2GRAY);
    equalizeHist(gray, gray);
    faceCascade.detectMultiScale(gray, faces, 1.1, 5, 0, Size(80,80));

    bool foundResident = false;
    for (auto& face : faces) {
        int label; double confidence;
        recognizer->predict(resized, label, confidence);

        if (label != -1 && confidence < RECOGNITION_THRESHOLD) {
            foundResident = true; // Cu dan hop le
            // Ve khung xanh + thong tin
        } else {
            // Ve khung do + canh bao Nguoi la
        }
    }

    // Chi gui lenh khi thang may dang o trang thai IDLE
    bool trigger = (foundResident && elevator.state == ELEVATOR_IDLE);
    if (updateElevator(elevator, trigger, bestResident, bestFaceCrop)) {
        // Gui 1 byte ma tang qua cong COM3 (chong spam)
        char floorChar = '0' + elevator.targetFloor;
        sendData(hSerial, floorChar);
    }

    drawElevatorUI(elevator); // Ve giao dien thang may
    imshow("Camera AI", frame);
    if (waitKey(30) == 'q') break;
}
\end{lstlisting}

\paragraph{Kết quả kiểm thử cuối tuần 2:}
\begin{table}[H]
    \centering
    \renewcommand{\arraystretch}{1.4}
    \begin{tabular}{|p{5cm}|p{3cm}|p{6.5cm}|}
        \hline
        \rowcolor{lightblue}
        \textbf{Chức năng kiểm thử} & \textbf{Kết quả} & \textbf{Ghi chú} \\
        \hline
        Đăng ký khuôn mặt 3 cư dân & \textcolor{green!60!black}{\bfseries Đạt} & Chụp đủ 30 mẫu/cư dân, lưu đúng cấu trúc thư mục. \\
        \hline
        Huấn luyện mô hình LBPH & \textcolor{green!60!black}{\bfseries Đạt} & Train $\approx$90 ảnh, lưu file \texttt{lbph\_model.yml}. \\
        \hline
        Nhận diện cư dân hợp lệ & \textcolor{green!60!black}{\bfseries Đạt} & Độ chính xác $>92\%$ (confidence $< 80$). \\
        \hline
        Phát hiện người lạ & \textcolor{green!60!black}{\bfseries Đạt} & Hiển thị cảnh báo đỏ, không kích hoạt thang máy. \\
        \hline
        Giao diện mô phỏng thang máy & \textcolor{green!60!black}{\bfseries Đạt} & Mở/đóng cửa mượt, di chuyển đúng tầng, $\approx$33 FPS. \\
        \hline
        Giao tiếp cổng COM3 & \textcolor{green!60!black}{\bfseries Đạt} & Gửi 1 byte mã tầng thành công; chế độ DEMO khi không có phần cứng. \\
        \hline
        Cache model (không train lại) & \textcolor{green!60!black}{\bfseries Đạt} & Khởi động lần 2 tải \texttt{.yml} tức thì, không huấn luyện lại. \\
        \hline
    \end{tabular}
    \caption{Bảng kết quả kiểm thử cuối Tuần 2}
\end{table}

\subsection{Tổng kết tuần 2}

\begin{table}[H]
    \centering
    \renewcommand{\arraystretch}{1.4}
    \begin{tabular}{|p{7cm}|p{7cm}|}
        \hline
        \rowcolor{lightblue}
        \textbf{Công việc đã hoàn thành} & \textbf{Ghi chú / Vấn đề phát sinh} \\
        \hline
        Cài đặt môi trường build MSYS2 + CMake thành công. & Gặp lỗi linker thiếu thư viện \texttt{-lopencv\_face}, đã sửa trong \texttt{CMakeLists.txt}. \\
        \hline
        Module DatabaseManager hoàn chỉnh (CRUD + quản lý ảnh). & Cần xử lý trường hợp tên cư dân trùng (UNIQUE constraint). \\
        \hline
        Nhận diện khuôn mặt LBPH hoạt động, cache model \texttt{.yml}. & Ngưỡng \texttt{RECOGNITION\_THRESHOLD = 80} cần tinh chỉnh theo ánh sáng thực tế. \\
        \hline
        State Machine thang máy 7 trạng thái hoạt động đúng. & Không có vấn đề phát sinh. \\
        \hline
        Giao diện OpenCV mô phỏng thang máy mượt mà. & Cần tối ưu khi nhiều khuôn mặt xuất hiện cùng lúc. \\
        \hline
        Tích hợp COM3, chạy demo end-to-end thành công. & Cổng COM10 trên máy mới cần thay tên cổng trong code. \\
        \hline
    \end{tabular}
    \caption{Tổng kết công việc Tuần 2}
\end{table}

\newpage

% =======================================================
% TỔNG KẾT CHUNG
% =======================================================
\section{TỔNG KẾT HAI TUẦN THỰC TẬP}

\subsection{Những gì đã đạt được}
\begin{itemize}[leftmargin=1.5cm]
    \item[$\checkmark$] \textbf{Tuần 1:} Nắm vững yêu cầu dự án, hoàn thành thiết kế hệ thống được Mentor phê duyệt.
    \item[$\checkmark$] \textbf{Tuần 2:} Triển khai thành công toàn bộ hệ thống SmartElevator Camera bằng C++17.
    \item[$\checkmark$] Áp dụng được kiến thức OOP, xử lý ảnh (Haar Cascade, LBPH), quản lý CSDL (SQLite3) và lập trình hệ thống (Win32 API, Serial COM) vào thực tế.
    \item[$\checkmark$] Sản phẩm hoạt động end-to-end: Camera AI $\rightarrow$ Nhận diện $\rightarrow$ Điều khiển thang máy mô phỏng.
\end{itemize}

\subsection{Kế hoạch các tuần tiếp theo}
\begin{enumerate}[leftmargin=1.5cm]
    \item \textbf{Tuần 3:} Tối ưu hóa độ chính xác nhận diện -- tinh chỉnh ngưỡng, thu thập thêm ảnh mẫu đa điều kiện ánh sáng.
    \item \textbf{Tuần 4:} Thử nghiệm với phần cứng thực tế (MCU), viết báo cáo kỹ thuật tổng hợp và chuẩn bị demo cho Mentor.
\end{enumerate}

\subsection{Nhận xét của sinh viên}

Hai tuần thực tập đầu tiên tại VNPT Media đã mang lại cho sinh viên những kinh nghiệm thực tiễn quý báu: từ quy trình làm việc chuyên nghiệp trong môi trường doanh nghiệp, cách phân tích và thiết kế hệ thống theo tài liệu SRS, đến việc tích hợp các công nghệ phức tạp (AI nhận diện ảnh, CSDL, giao tiếp phần cứng) vào một sản phẩm hoàn chỉnh bằng C++. Sinh viên sẽ tiếp tục phấn đấu hoàn thiện sản phẩm theo hướng dẫn của Mentor trong các tuần còn lại của kỳ thực tập.

\vspace{2.5cm}
\begin{table}[H]
    \centering
    \begin{tabular}{p{7cm}p{7cm}}
        \centering \textbf{Sinh viên thực tập} & \centering \textbf{Mentor xác nhận} \\[0.4cm]
        \centering \textit{(Ký và ghi rõ họ tên)} & \centering \textit{(Ký và ghi rõ họ tên)} \\[3cm]
        \centering \textbf{Nguyễn Đức Tuấn} & \centering \textbf{..............................} \\
    \end{tabular}
\end{table}

\end{document}
